% ============================================================================
\chapter{模曲线作为黎曼面}
\label{ch:riemann-surfaces}
% ============================================================================
% 对应 Diamond \& Shurman Chapter 2

在第一章，我们建立了商空间
$Y(\Gamma) = \Gamma \backslash \HH$ 和紧化模曲线
$X(\Gamma) = \Gamma \backslash \HH^*$。
但我们只说它是一个"商空间"——这个空间究竟有什么额外的数学结构？

答案是：$X(\Gamma)$ 是一个\textbf{紧黎曼面}（compact Riemann surface）。
黎曼面是最基础的一类复流形——每个点附近看起来像复平面的一小块。
正是这个复结构使我们能谈论"全纯函数"、"微分形式"和"亏格"，
从而把模形式解释为曲面上的微分形式，打通复分析与代数几何。

本章的主线如下：
\begin{itemize}
  \item \textbf{§\ref{sec:riemann-surface-def}} 介绍黎曼面的严格定义。
  \item \textbf{§\ref{sec:y-gamma-complex-structure}}
    把模曲线 $Y(\Gamma)$ 做成黎曼面：普通点与椭圆点各有不同的局部坐标。
  \item \textbf{§\ref{sec:compactification-X}}
    在尖点处添加局部坐标，得到紧黎曼面 $X(\Gamma)$。
  \item \textbf{§\ref{sec:riemann-hurwitz}}
    用 Riemann-Hurwitz 公式计算 $X(\Gamma)$ 的亏格。
  \item \textbf{§\ref{sec:modular-forms-differentials}}
    揭示模形式与全纯微分形式的对应关系。
\end{itemize}


% ============================================================================
\section{黎曼面：复平面的推广}
\label{sec:riemann-surface-def}
% ============================================================================

在复分析中，我们研究复平面 $\C$ 上的全纯函数。
黎曼面的想法是：把"局部像复平面"的性质变成全局对象的定义，
允许曲面在整体上有非平凡的拓扑。

先给一个直观图景：把一张纸卷成圆柱，每个点附近仍然"看起来像平面"，
但整体结构不同于平面。黎曼面就是这样——局部都像 $\C$，
但全局可以有洞、有亏格。

\begin{definition}[坐标卡与复图册]
\label{def:coordinate-chart}
\index{坐标卡 (coordinate chart)}\index{复图册 (complex atlas)}\index{黎曼面 (Riemann surface)}
设 $X$ 为 Hausdorff 拓扑空间。$X$ 上的一个\textbf{坐标卡}（coordinate chart）
是一对 $(U_\alpha, \varphi_\alpha)$，其中 $U_\alpha \subset X$ 是开集，
$\varphi_\alpha\colon U_\alpha \to V_\alpha \subset \C$ 是同胚（$V_\alpha$ 为开集）。

一族坐标卡 $\{(U_\alpha, \varphi_\alpha)\}$ 称为\textbf{复图册}（complex atlas），若：
\begin{enumerate}
  \item \textbf{覆盖}：$X = \bigcup_\alpha U_\alpha$；
  \item \textbf{全纯相容}：对任意两个坐标卡，转换映射
    \[
      \varphi_\beta \circ \varphi_\alpha^{-1}\colon
      \varphi_\alpha(U_\alpha \cap U_\beta) \to \varphi_\beta(U_\alpha \cap U_\beta)
    \]
    是\textbf{全纯函数}（biholomorphism）。
\end{enumerate}
带有极大复图册的连通 Hausdorff 空间 $X$ 称为\textbf{黎曼面}（Riemann surface）。
\end{definition}

% figures/ch02-riemann-surface-charts.tex
% 黎曼面的坐标卡与转换映射示意图
\begin{figure}[ht]
\centering
\begin{tikzpicture}[
  >=Stealth,
  every node/.style={font=\small},
]

% --- 黎曼面轮廓（用手工路径模拟椭圆曲面） ---
\draw[thick, fill=blue!8]
  plot[smooth cycle, tension=0.6]
  coordinates {(-2,0.8) (-0.5,1.5) (1.5,1.2) (2.2,0) (1.5,-1.2) (-0.3,-1.4) (-2,-0.5)};
\node at (0, 1.8) {$X$（黎曼面）};

% --- 坐标卡 U_alpha ---
\draw[thick, fill=orange!30, opacity=0.7]
  plot[smooth cycle, tension=0.7]
  coordinates {(-1.2,0.6) (-0.2,0.9) (0.3,0.2) (-0.2,-0.5) (-1.2,-0.3)};
\node at (-0.5, 0.2) {$U_\alpha$};

% --- 坐标卡 U_beta ---
\draw[thick, fill=cyan!30, opacity=0.7]
  plot[smooth cycle, tension=0.7]
  coordinates {(0.0,0.7) (1.0,1.0) (1.6,0.3) (1.1,-0.5) (0.1,-0.3)};
\node at (0.85, 0.25) {$U_\beta$};

% --- 交集标记 ---
\node[font=\footnotesize, text=gray!70!black] at (0.1, -0.9) {$U_\alpha \cap U_\beta$};
\draw[->, gray!60, shorten >=2pt] (0.1,-0.75) to[bend left=10] (0.35, 0.1);

% --- 左边 C 平面 (V_alpha) ---
\draw[fill=orange!10, draw=orange!60!black, thick, rounded corners=4pt]
  (-4.5,-3.2) rectangle (-2.5,-1.8);
\node at (-3.5, -2.5) {$V_\alpha \subset \C$};
\node[font=\footnotesize, text=gray] at (-3.5, -3.0) {$\varphi_\alpha(U_\alpha \cap U_\beta)$};
\draw[fill=orange!25, draw=orange!80!black, rounded corners=2pt]
  (-4.2,-2.9) rectangle (-3.3,-2.1);

% --- 右边 C 平面 (V_beta) ---
\draw[fill=cyan!10, draw=cyan!60!black, thick, rounded corners=4pt]
  (2.5,-3.2) rectangle (4.5,-1.8);
\node at (3.5, -2.5) {$V_\beta \subset \C$};
\node[font=\footnotesize, text=gray] at (3.5, -3.0) {$\varphi_\beta(U_\alpha \cap U_\beta)$};
\draw[fill=cyan!25, draw=cyan!80!black, rounded corners=2pt]
  (3.3,-2.9) rectangle (4.2,-2.1);

% --- 箭头：坐标映射 ---
\draw[->, thick, orange!70!black]
  (-0.5, -0.5) to[out=-120, in=60]
  node[left, pos=0.5] {$\varphi_\alpha$}
  (-3.5, -1.8);

\draw[->, thick, cyan!70!black]
  (0.85, -0.35) to[out=-60, in=120]
  node[right, pos=0.5] {$\varphi_\beta$}
  (3.5, -1.8);

% --- 转换映射箭头 (底部) ---
\draw[->, thick, purple!60!black, bend right=15]
  (-2.5, -2.5) to
  node[below, pos=0.5] {$\varphi_\beta \circ \varphi_\alpha^{-1}$（全纯双射）}
  (2.5, -2.5);

\end{tikzpicture}
\caption{黎曼面的坐标卡结构。曲面 $X$ 上的每个坐标卡 $(U_\alpha, \varphi_\alpha)$
把一小块区域同胚地映到 $\C$ 的开子集。两个坐标卡重叠时，
转换映射 $\varphi_\beta \circ \varphi_\alpha^{-1}$ 必须是全纯双射，
这保证了"全纯性"概念不依赖坐标卡的选择。}
\label{fig:riemann-surface-charts}
\end{figure}


\textbf{为什么需要"转换映射全纯"？}
两个坐标卡 $(U_\alpha, \varphi_\alpha)$ 和 $(U_\beta, \varphi_\beta)$
可以看成同一个点 $p \in U_\alpha \cap U_\beta$ 的两种"坐标系"。
在 $\varphi_\alpha$ 的坐标中，$p$ 对应 $z_\alpha = \varphi_\alpha(p) \in \C$；
在 $\varphi_\beta$ 的坐标中，$p$ 对应 $z_\beta = \varphi_\beta(p) \in \C$。
转换映射 $z_\beta = (\varphi_\beta \circ \varphi_\alpha^{-1})(z_\alpha)$ 全纯，
意味着"在不同坐标系中写出来的函数，全纯性不随坐标选择变化"。
这就是黎曼面上"全纯函数"概念有意义的根本原因。

\begin{example}[最简单的黎曼面：$\C$ 和开子集]
\label{ex:C-is-riemann-surface}
$\C$ 本身是黎曼面，图册只有一个坐标卡 $(U, \varphi) = (\C, \mathrm{Id})$。
更一般地，$\C$ 的任何连通开子集（如 $\HH$、单位圆盘 $\mathbb{D}$）也是黎曼面，
图册也只需一个坐标卡。
\end{example}

\begin{example}[黎曼球面 $\mathbb{P}^1(\C)$]
\label{ex:riemann-sphere}
\index{黎曼球面 (Riemann sphere)}\index{P1C@$\mathbb{P}^1(\mathbb{C})$}
$\mathbb{P}^1(\C) = \C \cup \{\infty\}$ 是最重要的紧黎曼面，拓扑上是一个球面。
定义两个坐标卡：
\begin{itemize}
  \item $U_0 = \C$，$\varphi_0(z) = z$（标准坐标）；
  \item $U_\infty = \C \setminus \{0\} \cup \{\infty\}$，
    $\varphi_\infty(z) = 1/z$（$\varphi_\infty(\infty) = 0$）。
\end{itemize}
在 $U_0 \cap U_\infty = \C \setminus \{0\}$ 上，转换映射为
$\varphi_\infty \circ \varphi_0^{-1}(z) = 1/z$，这是全纯函数。故 $\mathbb{P}^1(\C)$ 是黎曼面。

拓扑直觉：把复平面"卷起来"加上无穷远点，就得到一个球面
（球极投影，stereographic projection）。
\end{example}

\begin{example}[复环面是黎曼面]
\label{ex:torus-riemann-surface}
复环面 $\C/\Lambda$ 也是黎曼面。对任意 $p \in \C/\Lambda$，
取一个不包含格点的小邻域 $U$，
则投影 $\pi\colon \C \to \C/\Lambda$ 在 $\pi^{-1}(U)$ 的某个连通分支上是同胚，
其逆就给出坐标卡。转换映射是平移 $z \mapsto z + \omega$（$\omega \in \Lambda$），
显然全纯。

这与\cref{sec:complex-tori}中我们把 $\C/\Lambda$ 作为复 Lie 群的讨论一致。
\end{example}

\begin{definition}[黎曼面之间的全纯映射]
\label{def:holomorphic-map-riemann}
\index{全纯映射 (holomorphic map)}
设 $f\colon X \to Y$ 是黎曼面之间的连续映射。
$f$ 称为\textbf{全纯的}，若对 $X$ 和 $Y$ 上任何相容的坐标卡
$(U_\alpha, \varphi_\alpha)$ 和 $(V_\beta, \psi_\beta)$，
局部表达式
\[
  \psi_\beta \circ f \circ \varphi_\alpha^{-1}\colon
  \varphi_\alpha(U_\alpha \cap f^{-1}(V_\beta)) \to \C
\]
是全纯函数。

双全纯映射（biholomorphic map，即全纯且有全纯逆）也称为\textbf{双全纯同构}。
\end{definition}

\textbf{紧黎曼面上的全纯函数。}
若 $X$ 是紧黎曼面，则 $X$ 上没有非常数全纯函数。
这是因为：全纯函数的模在紧集上取到最大值，
由最大模原理，它必是常数。
这一事实说明我们应该研究紧黎曼面上的\emph{亚纯函数}，
而不是全纯函数——模形式（及其商 $j$-不变量等）正是这类对象。
